\section{The Apples of the Hesperides – Hercules’ Eleventh Labor}

We are nearing the end of the \emph{Dodekathalon}. We have seen how each of the Labors presupposes a certain internal struggle that is foundational to the soul, like a childbirth, since the Western path differs from archaic spirituality (eg., Vedic) in that it seeks to sublimate the material, rather than withdrawing away from it (although it is more accurate to state that the withdrawal becomes a foundational step that is transcended by the fullness of God' will which irradiates from the human heart, outward, and shines through matter). That is, a man must struggle with his passions, because they are bound up in an initial false self. We have seen Hercules pass through the stages inevitable to the reluctant hero: he has become popular with the masses as an outlaw-protector, he has made several mistakes (and recovered from them) which recapitulate the initial madness and sin which induced his quest, and he has discovered that enemies (eg., Hippolyta) are being converted to his side. He has begun to receive overt notice and aid from the celestial beings.

\textbf{Johannes Trimemius\footnote{\url{http://www.renaissanceastrology.com/trithemius.html}} writes that ``Study generates knowledge (cognition); knowledge prepares love; love, similarity; similarity, communion; communion, virtue; virtue, dignity; dignity, power; and power performs the miracle.''} This is the Christian sequence of Incarnational or natural magic. It is magic because God has ratified His oath from heaven in the Advent of the Son. It is natural, both because it was known through the \emph{Logos Tomeus} prior to historical event, and because, since the blessed Advent, Supernature is now overtly intervening and transforming, or uplifting, Nature. Einstein reminds us that ``Theory determines what we can observe'', so this is not a novel doctrine, but the re-discovery and energizing of what lay potential in Nature as seeds or signs of Logos, from the foundations of the Cosmos. Christianity is a religion of self-consciousness : that is, the Truth had always been there, but perceiving it clearly (with Calvin's spectacles) induces a change in the object of apperception, since Heaven and Earth began to ratify one another's oaths. This is the ``love-story'' of the Gospel.

Hercules' story is written in the stars, in the constellations, quite literally, but also symbolically, as each Labor corresponds to a progression of the Zodiac in the world of the cross. It is human, in that the figure is a warrior, an appropriate caste for the age of the Kali Yuga, in which man must struggle back and fight forward to bring about again the Age of Gold. It is also humane, in that Hercules stands for the classical Mediterranean World, the world of Man which gave birth to the Medieval Everyman, but which has roots in the Age of the Gods. The fact that this world decayed into that of animals (with the rise of Empiricism and Nominalism and Science) and that of machines (the Modern Era) and even that which is sub-human (our own time) should not dismay us: the subtle still rules the dense. Good has to interpenetrate that which is lowest. For even Evil cannot be created without first perverting that which is Good. So we see that, in the words of Mencius Moldbug, \emph{the Past is not a province of the Present, rather, the Present is a mere tiny slice and province (at best) of the Past.} So Evil (then) is merely a slice or province or truncated world of what is Good, what is God.

So we stand at a unique time: we have to ``go back'' to the antique past, because having collapsed into indeterminate quantum and sub-atomic spiritual conditions in the present, there is no way of emerging from this Chaos without re-instating a Traditional Order, and the only one that exists (and it still, physically even, exists) is in the past . To appeal to the Middle Ages is be novel in two senses: firstly, no one imagines it can be done, so that should be unique and interesting, and secondly, it is nearest to us, or ``newest'', in the sequence of Ages. The future exists indeterminately, within our collapsed state, and it is a privilege and rare opportunity to be of those who will ``create'' the Future by invoking the pattern of the Past. We are the keepers of the secret Fire, and the bearers of the apple seeds from the Hesperides, out of the deluge, into the Garden.

Hercules lived in such a time. Hera had come, and degeneracy reigned, because men had forgotten the reason for the sacrifices and mysteries. They performed them in ignorance, or as compromises with lower powers, since the process of entropy invariably proceeds by the dichotomy of a ``Two'' which does not understand the ``One'' that provides it birth and Order. Look at how men think today. Do they not invariably pose false dichotomies, and then flatten out all meaning and nuance in order to make the supposed synthesis plausible? This is because the power that opposes man's true destiny, and which feeds off of it (we can call It and Them ``Moriarty'' if we like) and requires man's ignorance in order to continue to drag out its existence, has access to men's Mind (collective and individual) through the disorder of passion and the perversion of Reason to slave to that Passion. And this is why Hercules is ``at war'', taking heaven by storm, within and without. Hercules fights, initially, at times and places not of his own choosing, although even there we see a tiny sliver of man's free will which allows some significant freedom and decision-making (after all, he did strangle a serpent in his crib), but as the process matures magically, Hercules begins to enjoy the fruits of conquest, and to make war at times and places more of his own choosing.

Man is growing up. We see this in this episode, because it is virtually an entire twelve labors, in and of itself. Hercules doesn't even know where the Apples are, and has to wrestle the Old-Man-of-the-Sea, who keeps changing shapes while being held. Certainly, this has something to do with achieving a meditative state of mind, with stilling the flux, until the image emerges in the still pool, as Tomberg readers will recognize. He then has to slay Antaeus, by holding him aloft so that he does not touch the earth, and crushing him. Again, this can be related to asceticism, which involves severing the ``natural'' drive of desire from its magnetized source in the passions, to make it serve something not purely of this world. Journeying to Egypt, he is imprisoned by a king (again we see conspiracy and political conflict), who desires to sacrifice Hercules himself for the yearly offering. Indeed, this is done routinely in the Brave New World. It is the liberal version of ``Knockout Game'', in which some endangered species of latent heretic is dragged out of his abode and thrashed around for vicarious thrill and for the pure pleasure of witch-hunting\footnote{\url{http://unqualified-reservations.blogspot.com/2013/09/technology-communism-and-brown-scare.html}}. Have you tried this lately? It's lots of fun. And you don't even have to have any real courage, because, after all, the people you are attacking have no actual power to hit back. Plus, you can feel like the Pharisees did when they tried to stone the woman they had committed adultery with. But, really, it's just a species of perverted and ignorant dark magic, because the \emph{faux Liberaux} deep down just think that every witch they burn will help ensure the pleasure of the Multi-Cult Gods. Don't burn too many: one or two a year will do. After all, who will do the work around this insane asylum?

After this hold up, Hercules manages to break through. We don't hear of too many accidental killings this time around. Hercules is just thinking about those apples, and apparently, even his enemies are becoming less inclined to use open violence against him. When he arrives at the Hesperides, he finds that he himself cannot pick the Apples. In an interesting twist, which is of worthy note, he tricks Atlas first into doing it for him (while he holds up the world on his back) \& then lies to Atlas that he needs to retrieve his cloak to make things easier for holding the globe, will Atlas hold it up momentarily?

The Lesson here is that, apparently, the hero is no longer bound by a conventional code of honor. Like Odysseus, the cunning one of many stratagems, Hercules is allowed (apparently) to use help (this was disqualified, as you remember, in the case of the Augean Stables and the Hydra), and to use it dishonestly. Apparently, the Apples are a prize which does not require perfect virtue or honor to retrieve.

As a Christian, I would like to re-interpret this passage, if it were permissible, but honesty forbids it. Here we reach into the mystery of Good \& Evil, where the law or ban of the true sovereign reaches a point that places that entity beyond what is commonly conceived as right and wrong. To attain some perspective, and try to square the Circle, we will note that even the decent, salt-of-the-earth common man lives in some degree of tension and hypocrisy\footnote{\url{http://www.atimes.com/atimes/Front\_Page/GA19Aa02.html}}, when you scratch him deep enough. After all, we all shop at Wal-Mart.

In this instance, I will say that Hercules returned at a later date, building the pillars of Hercules, belatedly, thus liberating Atlas much as he freed Prometheus. But he still lied. Indeed, it was a lie extorted in bad faith by Atlas, as the apples would do Hercules little good, were he to stand there forever hoisting the world orb. Mouravieff deals with some of this in \emph{Gnosis}. The lie can never be to the Self, although Life forces many lies upon us, in many contexts, in which (at the least) we allow people to see what we perceive them to be mis-perceiving, without correcting them, because to ``always tell the truth'' in a bald manner at all times would ensure that we spent our lives much like Don Quixote spent his. I might also add, that the spiritual laws of Fairie Land, where the Apples grow, are not like unto the laws of Middle Earth, where the common code of the Hobbit must tend to prevail. Cunning is always a quality heavily valued in the Fairy Tale; we have to learn, as humans, that we are not always speaking to another person, but to an archetype, especially in a quest.

\emph{And what of ``let your Yes be Yes, and your No, No''?} As a Christian, I say that I am not yet the Son of God, and furthermore, will never be so in the same way that Christ succeeded in being so, as He never knew sin. The religion of the apostles is slightly different in modulation and form from the religion practiced by Christ Himself. This ``gap'' is not a proof that one is not understanding the Logos, but rather, a testimony to the innately flawed and tragic situation man has cognized, and seeks to transcend. Christ, after the Resurrection, was not the same Christ Who was laid in the tomb. If we see the Church after Christ's Advent, the promise is made to reconcile even those who are in Purgatory, those who are \emph{idiotes} or imperfect, those who dimly cling to the shreds of Christ's garment, can find salvation in the precincts of the outer Temple. This was the truth which Luther really could not quite come to terms with, and which rejection, incidentally, still marks his spiritual children in the modern world. For them, a single imperfection invalidates the entire corpus. This ensures that one will gravitate towards Alfred Jarry's dictum: \emph{Until we've destroyed everything, we've done nothing}. The Perfect becomes the Enemy of the Good.

The Christian path recognizes the Evil and the hypocrisy, but (as noted), the goal is not withdrawal from matter, but rather a shining through and sublimation of that matter, however, imperfectly, which transforms even Evil into Good. This is a delicate position, as it can come perilously close to ``Do Evil that Good may come'' or the doctrine of Progressivism that any broken eggs are worth the omelette. What I am trying to say here is something capable of distortion. It is rather that within us, it sometimes happens, when the hero is pushed to the white-hot extremities of contradiction during his quest, he will sometimes find himself acting (unconsciously perhaps, or consciously, bearing the guilt truly) with elements of his being that are not perfectly devoted to God. To give a concrete example, suppose I am awakened in the middle of the night by the noise of strangers entering my home. Now, as a good Christian, my first thought should be to ensure the safety, not only of my family, but even of those who have knocked in the door, if possible. Indeed, a pure and true Christian could square the circle, with the help of angels, by reading their thoughts, speaking words of wisdom, or even laying down his life to save both family and foe. However, what would most likely happen (because I know myself, especially at two o'clock in the morning) would be that a weapon would be procured and used. My greatest fear (being imperfect) would not be to fail to be a perfect Christian, since that is a given, but to allow harm to come to my family through personal cowardice \emph{or weakness or hesitation. In such a situation, it is likely that violence would occur.}

This is the mystery of the temple of Rimmon. It is analogous to the Earps deciding to clean up Tombstone themselves, rather than allow Sheriff Beehan to arrest them for disorderly conduct in the middle of a flagrant murdering spree perpetrated by those who own the Sheriff. In fact, almost every human being has experienced a situation, on some scale, of the same kind, in which to ``do right'' is actually to commit an even graver sin. What is wrong with modern Liberalism is not their recognition (they do this too) that sometimes you have to go against your kin in order to be Just to the stranger and the alien, but rather, the fact that they literally worship these strangers and outcasts as God Himself. And the fact that they do not allow such niceness of conscience or subtlety of distinction to be used by anyone else but those who also, like them, worship at the altar of Revolution and Degeneracy.

In this case, Hercules knows he has to seize the apples, and they are his provided he is willing to answer for the consequences -- he can return later and redeem his pledge. Hasn't he committed worse sins than this, and been forgiven? What is worse than slaughtering your wife and children in a fit of madness, or a woman who loves you as she brings your prized object? We see that Hercules is simply ``unequal'' to ordinary standards of right and wrong. This is completely different (as far as I can tell) from the Nietzschean idea that one must go beyond them entirely. In a word, Hercules is a gentleman, but with a bit of a rogue left in him. Indeed, anyone who seeks individual salvation will find that they become a gentleman, with a bit of something else streaked through, which remains (even in the redeemed state, in Heaven) as something like the silver lining of a storm cloud.

Like Cromwell\footnote{\url{http://www.youtube.com/watch?v=dRFrXGkYjhM}}, if you are willing to stand tall before God and answer for the blood on your gauntlets, then ``you may proceed''. \emph{Take what you want, and pay God for it}. But, had Cromwell subdued himself more thoroughly, and acted more in accordance with his caste, had he seen himself caught in the currents and dragged by Fate as a slave to the rising order of Revolution, would he have acted so?

With the ancients, we affirm that no man consciously would will Evil (at least before he was thoroughly corrupted, as this Life is a second chance, and perhaps not even then). Which is why it is the duty of every man of Order to subdue himself, that when he creates, he may create according to the whole Eye of Light, perceiving the Logos not through a glass darkly, but almost face-to-face.


\flrightit{Posted on 2013-12-21 by Logres}

\begin{center}* * *\end{center}

\begin{footnotesize}
\begin{sffamily}
\texttt{magicalcharms on 2013-12-26 at 12:22 said:}

thought-provoking, helpful and instructive thank you.


\hfill

\end{sffamily}
\end{footnotesize}

